% Generated by scripts/assemble.py; corrections belong in source/corrections.json.
\examhead{84}{哈尔滨工程大学 2024 年数学分析试卷}{83}
% p084-b01; source PDF page 84
\par\Needspace{6\baselineskip}
\blockanchor{p084-b01}
\question{1}令 $A, B$ 为任意的非空有界实数集, 定义\par
\nopagebreak[4]
% p084-b02; source PDF page 84
\blockanchor{p084-b02}
\displayformula{A + B = \{z \mid z = x + y, x \in A, y \in B \}, A - B = \{z \mid z = x - y, x \in A, y \in B \}.}
% p084-b03; source PDF page 84
\blockanchor{p084-b03}
证明: $\sup (A + B) = \sup A + \sup B$, $\sup (A - B) = \sup A - \inf B$.\par
% p084-b04; source PDF page 84
\blockanchor{p084-b04}
\question{2}计算 $\lim_{n\to \infty}\sin (n!\alpha \pi)$, 其中 $\alpha$ 为有理数.\par
% p084-b05; source PDF page 84
\blockanchor{p084-b05}
\question{3}求曲线 $\left(\frac{x^2}{a^2} + \frac{y^2}{b^2}\right)^2 = x^2 + y^2$ 所围平面图形的面积.\par
% p084-b06; source PDF page 84
\blockanchor{p084-b06}
\question{4}设  $f(x,y)$  在  $P(x_{0},y_{0})$  处可微,  $l_{1},l_{2},\cdots,l_{n}$  为  $P(x_{0},y_{0})$  点处给定的 n 个方向不同的单位向量, 两相邻向量夹角为  $\frac{2\pi}{n}$ , 计算  $\sum_{i=1}^{n}\frac{\partial f(x_{0},y_{0})}{\partial l_{i}}$ , 其中  $\frac{\partial f(x_{0},y_{0})}{\partial l_{i}}$  表示  $f(x,y)$  在  $P(x_{0},y_{0})$  处沿  $l_{i}$  的方向导数.\par
% p084-b07; source PDF page 84
\blockanchor{p084-b07}
\question{5}令 $x = \frac{1}{n} \sum_{k=1}^{n} x_k$, 其中 $x_i \in (0, \pi)$, 证明: $\prod_{k=1}^{n} \frac{\sin x_k}{x_k} \leq \left(\frac{\sin x}{x}\right)^n$.\par
% p084-b08; source PDF page 84
\blockanchor{p084-b08}
\question{6}求 $f(x) = [x] \sin \pi x$ 的单侧导数.\par
% p084-b09; source PDF page 84
\par\Needspace{6\baselineskip}
\blockanchor{p084-b09}
\question{7}如果 $f(x)$ 在 $[0,1]$ 上二次可微, $f''(x)$ 有界且黎曼可积. 证明:\par
\nopagebreak[4]
% p084-b10; source PDF page 84
\blockanchor{p084-b10}
\displayformula{\lim _ {n \rightarrow \infty} n ^ {2} \left(\int_ {0} ^ {1} f (x) \mathrm{d} x - \frac {1}{n} \sum_ {i = 1} ^ {n} f \left(\frac {2 i - 1}{2 n}\right)\right) = \frac {f ^ {\prime} (1) - f ^ {\prime} (0)}{24}.}
% p084-b11; source PDF page 84
\par\Needspace{6\baselineskip}
\blockanchor{p084-b11}
\question{8}令\par
\nopagebreak[4]
% p084-b12; source PDF page 84
\blockanchor{p084-b12}
\displayformula{u _ {n} (x) = \left\{ \begin{array}{l l} - n, & x \leq - n; \\ x, & - n <   x \leq n; \\ n, & x > n. \end{array} \right.}
% p084-b13; source PDF page 84
\blockanchor{p084-b13}
$f(x)$ 为实值函数, 证明: $f(x)$ 连续等价于对任意固定的正整数 $n$, $g_{n}(x) = u_{n}(f(x))$ 都是 $x$ 的连续函数.\par
% p084-b14; source PDF page 84
\blockanchor{p084-b14}
\question{9}假设 $f(x) \in C^{\infty}([0,1])$ 且满足下列条件:\par
% p084-b15; source PDF page 84
\blockanchor{p084-b15}
(1) $f(x)\neq 0$\par
% p084-b16; source PDF page 84
\blockanchor{p084-b16}
(2)  $f^{(n)}(x)=0,\ n=0,1,2,\cdots$\par
% p084-b17; source PDF page 84
\blockanchor{p084-b17}
(3) 实数序列 $\{a_{n}\}$ 满足级数 $\sum_{n=1}^{\infty} a_{n} f^{(n)}(x)$ 在 $[0,1]$ 上一致收敛.\par
% p084-b18; source PDF page 84
\blockanchor{p084-b18}
证明： $\lim_{n\to \infty}n!a_n = 0.$\par
% p084-b19; source PDF page 84
\blockanchor{p084-b19}
\question{10}证明: $H_{n}(x) = (-1)^{n} e^{x^{2}} \frac{\mathrm{d}^{n}}{\mathrm{d} x^{n}} (e^{-x^{2}})$ 所有根都是实数.\par
% p084-b20; source PDF page 84
\par\Needspace{6\baselineskip}
\blockanchor{p084-b20}
\question{11}计算第二类曲面积分\par
\nopagebreak[4]
% p084-b21; source PDF page 84
\blockanchor{p084-b21}
\displayformula{\iint_ {S} (f (x, y, z) + x) \mathrm{d} y \mathrm{d} z + 2 f (x, y, z) \mathrm{d} z \mathrm{d} x + (f (x, y, z) + z) \mathrm{d} x \mathrm{d} y.}
% p084-b22; source PDF page 84
\blockanchor{p084-b22}
其中 $f(x,y,z)$ 为连续函数，$S$ 为平面 $x - y + z = 1$ 在第四卦限的上侧.\par
% p084-b23; source PDF page 84
\blockanchor{p084-b23}
\question{12}设 $f(x, y)$ 定义在 $x^2 + y^2 \leq 1$，且具有连续的偏导数，$|f(x, y)| \leq 1$. 证明：在 $x^2 + y^2 \leq 1$ 内有一点 $(x_0, y_0)$，使得 $|f_x(x_0, y_0)|^2 + |f_y(x_0, y_0)|^2 < 16$.\par
